%% Die TeXnische Komoedie (dtk) article starter — DANTE e.V.'s journal.
%%
%% The dtk class itself is NOT included in this download: it ships with the
%% pinned compiler's TeX Live, so bundling a copy would only risk pinning you
%% to an older version than the compiler already has.
%%
%% REQUIRES LuaLaTeX: the class refuses to load under pdflatex/xelatex.
%%
%% BIBLIOGRAPHY: the class hardcodes \RequirePackage[backend=biber]{biblatex},
%% so this starter's own \addbibresource/\printbibliography drive biber (not
%% classic bibtex) via the compiler's biblatex/biber cycle.
%%
%% Class documentation: texdoc dtk   |   https://ctan.org/pkg/dtk
%% Licence: dtk is LPPL-1.3c, (c) Rolf Niepraschk, Herbert Voss.
\documentclass[ngerman]{dtk}
\usepackage{booktabs}
\addbibresource{main.bib}

\begin{document}
\title{Ein Leitfaden zum Setzen von Algorithmen mit der \texttt{dtk}-Klasse}
\Author{Lovelace}{Ada}%
    {Institut für Analytische Maschinen\\
     \Email{ada@example.edu}}
\Author{Turing}{Alan}%
    {Fakultät für Rechnende Maschinen\\
     \Email{alan@example.edu}}
\maketitle

\section{Einleitung}
Dieser Beispielartikel zeigt, wie ein Manuskript für \journalname{} mit der
Klasse \texttt{dtk} vorbereitet wird. Ersetzen Sie diesen Text durch Ihre
eigene Einleitung.

\section{Aufbau des Artikels}
Ein typischer Artikel gliedert sich in Abschnitte wie diesen. Absätze werden,
wie in der KOMA-Script-Tradition üblich, ohne Einzug und mit vertikalem
Abstand gesetzt.

\begin{itemize}
\item Listen werden mit Punkten eingeleitet.
\item Verschachtelung ist bis zu vier Ebenen tief möglich.
\end{itemize}

\section{Quelltext-Listings}
Die Klasse bindet das Paket \texttt{listings} bereits ein. Ein nummeriertes
Listing:

\begin{lstlisting}[style=number]
\begin{itemize}
\item Erster Punkt
\item Zweiter Punkt
\end{itemize}
\end{lstlisting}

Und ein Listing ohne Nummerierung:

\begin{lstlisting}[style=noNumber]
\begin{description}
\item[Docstrip] zerlegt eine .dtx-Datei in ihre Bestandteile.
\end{description}
\end{lstlisting}

\section{Tabellen}
\begin{table}[htbp]
  \caption{Ein einfaches Beispiel.}
  \label{tab:beispiel}
  \begin{tabular}{lrr}
    \toprule
    Bedingung & Anzahl & Anteil \\
    \midrule
    Basis     & 120    & 0{,}42 \\
    Neu       & 118    & 0{,}71 \\
    \bottomrule
  \end{tabular}
\end{table}

Tabelle~\ref{tab:beispiel} zeigt die Werte.

\section{Zusammenfassung}
Fassen Sie hier die wichtigsten Aussagen Ihres Artikels zusammen. Ein Zitat
\autocite{knuth1984texbook} zeigt, wie Literaturverweise mit der
\texttt{biblatex}-Schnittstelle der Klasse gesetzt werden.

\printbibliography

\end{document}
